% !TeX root = main.tex
% Judul Bab II (Pusat)
\begin{center}
	\setcounter{huruf}{0}
	\textbf{BAB II} \\
	\textbf{PEMBAHASAN}
\end{center}
\addcontentsline{toc}{section}{BAB II PEMBAHASAN}

\vspace{0.5cm}

% Menggunakan command \hurufsection untuk poin pembahasan pertama
\hurufsection{Pengertian Teknologi Informasi dalam Perpustakaan}{
	Teknologi informasi dalam konteks perpustakaan mencakup segala bentuk teknologi yang digunakan untuk memproses, menyimpan, dan menyebarkan informasi. Hal ini meliputi penggunaan perangkat keras, perangkat lunak, serta jaringan komunikasi yang mendukung operasional perpustakaan secara digital.\footcite{basuki2006} Secara fundamental, teknologi informasi di perpustakaan bukan sekadar alat bantu elektronik, melainkan satu kesatuan sistem yang mengintegrasikan pengelolaan data koleksi dengan layanan pengguna secara real-time.\footcite{basuki2006}

	Implementasi teknologi ini bertujuan untuk meningkatkan efisiensi layanan dan memudahkan pengguna dalam mengakses koleksi yang tersedia, baik dalam bentuk fisik maupun digital. Melalui sistem berbasis komputer, tugas-tugas rutin pustakawan seperti sirkulasi dan pengatalogan dapat diselesaikan dengan lebih akurat dan cepat, sehingga meminimalisir kesalahan manusia yang sering terjadi pada sistem manual.

	Lebih jauh lagi, pemanfaatan teknologi informasi memungkinkan perpustakaan untuk memperluas jangkauan layanannya melampaui batas fisik gedung. Perpustakaan kini mampu menyediakan akses informasi tanpa henti (24/7) melalui katalog daring dan pangkalan data digital.\footcite{darmono2007} Hal ini menciptakan ekosistem informasi yang lebih demokratis, di mana pemustaka dapat menemukan referensi yang dibutuhkan dengan hanya beberapa klik, sekaligus memungkinkan perpustakaan untuk menjalankan fungsi pelestarian digital terhadap naskah-naskah langka agar tetap terjaga kualitas informasinya di masa depan.
}

% Poin pembahasan kedua
\hurufsection{Penerapan Sistem Otomasi Perpustakaan}{
	Sistem otomasi perpustakaan merupakan aplikasi teknologi informasi yang digunakan untuk mengelola fungsi-fungsi utama perpustakaan secara terintegrasi, mulai dari pengadaan, pengatalogan, sirkulasi, hingga pengelolaan anggota. Implementasi sistem ini memungkinkan seluruh data bibliografi tersimpan dalam pangkalan data terpusat, sehingga proses temu balik informasi dapat dilakukan secara instan. Menurut Abdul Rahman Saleh, otomasi bukan sekadar memindahkan catatan manual ke komputer, melainkan upaya mengintegrasikan seluruh unit kerja perpustakaan ke dalam satu alur kerja digital yang koheren.\footcite{saleh2014otomasi}

	Dengan adanya otomasi, pustakawan dapat bekerja lebih cepat dan akurat, terutama dalam mengelola transaksi peminjaman dan pengembalian koleksi yang memiliki volume tinggi. Penggunaan teknologi kode batang (\textit{barcode}) atau \textit{Radio Frequency Identification} (RFID) dalam sistem otomasi secara signifikan mengurangi risiko kesalahan input data dan mempercepat proses inventarisasi koleksi. Wahyu Supriyanto menekankan bahwa efisiensi yang dihasilkan dari sistem otomasi memberikan ruang bagi pustakawan untuk lebih fokus pada pengembangan layanan literasi dan pendampingan pemustaka, daripada terjebak dalam tugas-tugas administratif yang repetitif.\footcite{supriyanto2008}

	Bagi pemustaka, manfaat utama dari sistem otomasi adalah kemudahan akses melalui \textit{Online Public Access Catalog} (OPAC). Layanan ini memungkinkan pencarian koleksi dilakukan tanpa batasan ruang dan waktu, serta memberikan informasi \textit{real-time} mengenai ketersediaan buku di rak. Dalam perkembangannya, sistem otomasi perpustakaan modern kini banyak yang bersifat \textit{open source}, seperti SLiMS (\textit{Senayan Library Management System}), yang memungkinkan perpustakaan kecil hingga menengah untuk melakukan digitalisasi layanan dengan biaya yang lebih terjangkau namun tetap memiliki fitur yang komprehensif.\footcite{mulyani2020}
}

% Poin pembahasan ketiga
\hurufsection{Tantangan Digitalisasi Informasi}{
	Meskipun teknologi memberikan banyak kemudahan, proses digitalisasi informasi juga menghadapi berbagai tantangan yang kompleks. Salah satunya adalah masalah hak cipta (\textit{copyright}) yang sering kali menjadi penghambat dalam penyebaran koleksi secara daring. Perpustakaan harus memastikan bahwa proses konversi dari dokumen fisik ke format digital tidak melanggar undang-undang hak cipta, terutama terkait dengan penggandaan dan distribusi karya intelektual. Menurut Putu Laxman Pendit, tantangan hukum ini menuntut pustakawan untuk tidak hanya menguasai aspek teknis, tetapi juga memahami etika dan regulasi mengenai kekayaan intelektual di ranah digital.\footcite{pendit2007}

	Selain aspek hukum, kebutuhan akan infrastruktur yang memadai dan pemeliharaan sistem secara berkala (\textit{system maintenance}) menjadi syarat mutlak agar data tetap aman dan dapat diakses dalam jangka panjang. Investasi pada perangkat keras (\textit{hardware}) seperti peladen (\textit{server}) berkapasitas tinggi dan sistem pencadangan (\textit{backup system}) yang mumpuni sangat diperlukan untuk mencegah kehilangan data akibat kegagalan sistem. Romi Satria Wahono menyatakan bahwa keterbatasan anggaran sering kali menjadi kendala bagi perpustakaan dalam melakukan pembaruan infrastruktur TIK, padahal teknologi terus berkembang dengan sangat cepat (\textit{high obsolescence rate}).\footcite{wahono2006}

	Yang tidak kalah krusial adalah masalah keamanan informasi (\textit{information security}) dan kerentanan terhadap serangan siber (\textit{cyber attacks}). Penjahat siber sering kali menargetkan pangkalan data perpustakaan untuk mencuri data anggota atau merusak integritas koleksi digital melalui \textit{malware} atau \textit{ransomware}. Oleh karena itu, penerapan protokol keamanan seperti enkripsi data dan \textit{firewall} yang kuat harus menjadi prioritas dalam pengelolaan perpustakaan modern. Di samping itu, masalah pelestarian digital (\textit{digital preservation}) juga menjadi perhatian penting, di mana format file digital harus terus diperbarui agar kedepannya tetap bisa terbaca oleh perangkat lunak.\footcite{subrata2014}
}